% 08-codegen.tex
\section{代码生成与构建}
\label{sec:codegen}

代码生成是唯一可以保持 479 个 \ac{op} 模式、187 个后向公式、13 个调度键、两个语言表面和 6 个构建工件相互一致而不会出现偏差的机制。合约是两个YAML文件；该引擎是一个由本地模式模型驱动的纯Python生成器；输出是十个分段标头和源代码，它们被编译为普通翻译单元。生成后没有任何内容是手工编辑的：每个工件的第一行读取 \texttt{main.py -- DO NOT EDIT} 并且如果生成器过时，构建将拒绝链接。

\subsection{合约：两个 YAML 文件}

运算符注册表 \texttt{native\_functions.yaml} 有 9\,445 行，当扩展重载和 \texttt{.out} 变体时，在 2\,809 \texttt{- func:} 记录中声明 479 个不同的运算符名称。每个条目绑定五个字段：\texttt{func}（FunctionSchema 字符串）、\texttt{variants}（\texttt{function} 和/或 \texttt{method}）、\texttt{dispatch}（从 \texttt{CPU}、\texttt{CUDA}、\texttt{Composite} 到内核符号的映射）、可选 \texttt{python\_defaults} 和可选 \texttt{device\_check} / \texttt{structured\_delegate} / \texttt{tags}。架构字符串是 C++ 签名、Python 绑定和调度程序密钥表的单一来源。变体决定面：\texttt{function} 成为静态成员 \texttt{Tensor::zeros} 和模块函数 \texttt{tensorplay.zeros}； \texttt{method} 成为实例成员 \texttt{tensor.view}。代表性工厂条目会内联拼写其默认值，例如\ \texttt{func: zeros(int[] size, *, ScalarType dtype=Float32, Device device=cpu, bool requires\_grad=False) -> Tensor} 与 \texttt{variants: function} 和 \texttt{dispatch: \{CPU: zeros\_cpu, CUDA: zeros\_cuda\}}。就地条目具有以下类型的可变性：\texttt{add\_(Tensor(a!) self, Tensor other, Scalar alpha=1) -> Tensor(a!)} 和 \texttt{variants: method}。

为 1\,330 行（标签处为 1\,324 行），其中包含 345 条 \texttt{- name:} 记录，按基本名称分组后可折叠为 187 个向后公式。每个记录都以符号方式命名前向参数梯度： \texttt{name: add(Tensor self, Tensor other, Scalar alpha=1) -> Tensor} 、 \texttt{self: grad} 和 \texttt{other: grad * alpha} 等。\ YAML 从不携带 C++ 类型或调度键；这些是从它引用的本机条目推断出来的。公式可以调用分派的操作 (\texttt{mul}, \texttt{sum\_to\_shape}) 和由为向后节点发射器提供动力的相同表达式 AST 解析的自由函数。随机操作会自动添加 \texttt{nondeterministic\_seeded} 标签，并且 \texttt{Vulkan} 密钥在 CPU 版本上保留保留。

除了自己的 \texttt{*.py} 源之外，这两个文件是 Python 管道允许读取的唯一输入。从结构上检查不变量： \texttt{parse\_native\_yaml} 通过本地模式管道解析同一文件两次，将 \texttt{int64\_t} 规范化为 \texttt{int} ，将 \texttt{Composite} 规范化为 \texttt{CompositeExplicitAutograd} ，并在任何删除的记录、分派投影不匹配、变体不匹配或运算符名称冲突时引发 \texttt{ValueError} 。

\begin{tpnote}[Note: YAML Is the Single Source of Truth---Everything Else Is Derived]
每个生成的表面都是两个 YAML 文件的纯函数。 C++ 签名、Python 绑定、\texttt{.pyi} 存根、调度程序密钥表、向后节点和自动转换注册都是 \emph{derived artifacts}：从相同的 YAML 重新生成始终会重现字节相同的输出，并且 \texttt{DO NOT EDIT} 标头使手动修改可检测到，而不仅仅是阻止。贡献者的实际结果是： \emph{never} 编辑生成的文件以修复编译错误——该修复属于 YAML 或生成器，并将覆盖下一个构建时的手动编辑。单一来源规则使完整性检查变得有意义：它们验证从 YAML 到工件的 \emph{projection} 是完整的和单射的，这只是一个明确定义的声明，因为方向是单向的。
\end{tpnote}

\begin{lstlisting}[caption={and fragments.}, label=lst:yamlfrag, style=pystyle]
- func: add.Tensor(Tensor self, Tensor other, *, Scalar alpha=1) -> Tensor
  variants: function, method
  dispatch: {CPU: add_cpu, CUDA: add_cuda}
  python_defaults: {alpha: 1}
- func: add_.Tensor(Tensor(a!) self, Tensor other, *, Scalar alpha=1) -> Tensor(a!)
  variants: method
  dispatch: {CPU: add__cpu, CUDA: add__cuda}
  device_check: NoCheck   # lists checked element-wise elsewhere
- name: add(Tensor self, Tensor other, Scalar alpha=1) -> Tensor
  self: grad * alpha
  other: grad * alpha
\end{lstlisting}

\subsubsection{YAML 合同完整性检查：\texttt{parse\_native\_yaml}}

\texttt{parse\_native\_yaml} 是 YAML 文本和每个生成器之间的单一门户。它执行七项结构检查，共同保证合同的完整性和单射性。该实现加载文件两次：一次通过本地模型的 \texttt{YamlLoader}，一次通过独立的结构验证路径。没有模式可以绕过第二次解析。

\begin{table}[H]
\centering
\caption{完整性检查和--\texttt{752}。}
\label{tab:yaml_checks}
\small
\begin{tabular}{lll}
\toprule
查看 & Site & 失败 \\
\midrule
删除记录 &  & \texttt{ValueError: schema parser dropped records} if \texttt{len(parsed) < len(prepared)} \\
重复运算符 & , \texttt{1026} & \texttt{duplicate parsed operator schema} / \texttt{duplicate operator schema \{func\_name\}} \\
缺失/额外 &  & \texttt{record set mismatch: missing=\dots extra=\dots} \\
模式投影 &  & \texttt{schema projection mismatch: \{actual\} != \{预期\}} \\
模式种类 &  & \texttt{mutable} \\
变体集 &  & \texttt{variant mismatch: \{actual\} != \{预期\}} \\
调度投影 &  & \texttt{dispatch projection mismatch: \{actual\} != \{预期\}} \\
\bottomrule
\end{tabular}
\end{table}

具体来说，\texttt{\_prepare\_schema\_entry} 在验证之前规范化每个项目：它剥离 \texttt{python\_defaults}，将 \texttt{int64\_t} 重写为 \texttt{int}，将 \texttt{Composite} 提升为 \texttt{CompositeExplicitAutograd}，为名称包含 \texttt{rand} 或参数包含 \texttt{dropout} ( \texttt{\_requires\_seeded\_tag}) 的任何模式注入 \texttt{nondeterministic\_seeded}，并为工厂合成 \texttt{CompositeExplicitAutograd} 调度 (\texttt{new\_*} / \texttt{*\_like} / 仅 TensorOptions）通过 \texttt{cpp.name}。 \texttt{\_validate\_dispatch\_projection} 然后根据 \texttt{backend\_metadata} 验证 YAML \texttt{dispatch} 映射；仅忽略 \texttt{\_\_line\_\_}。 \texttt{\_validate\_native\_projection} 为本地 \texttt{NativeFunction} 断言四个相等：字符串化 \texttt{FunctionSchema}、\texttt{SchemaKind} (\texttt{functional} / \texttt{inplace} / \texttt{mutable})、变体名称集和 \texttt{cpp\_no\_default\_args} 集。顶级 \texttt{parse\_native\_yaml} 循环通过 \texttt{seen\_schemas} 删除精确重复的 \texttt{- func:} 字符串（合并工件），但将解析为相同 \texttt{func\_name} （过载冲突）的任何两个不同字符串视为错误。最后，设置相等可确保没有名称丢失或多余。

\begin{lstlisting}[caption={Completeness check diagnostics (abridged from --\texttt{1044}).}, label=lst:yamlchecks, style=pystyle]
# model.py:707  dispatch mismatch
raise ValueError(f"dispatch projection mismatch for {f.func_name}: "
                 f"{actual} != {expected}")
# model.py:982  dropped records
raise ValueError(f"schema parser dropped records for {path}: "
                 f"{len(prepared)} > {len(parsed.native_functions)}")
# model.py:1038 record set mismatch
if configured_names != reference_names:
    raise ValueError(f"schema parser record set mismatch: "
                     f"missing={missing[:8]} extra={extra[:8]}")
\end{lstlisting}

\subsection{架构模型：（1\,055 行）}

模型层是本地模式记录的薄投影。 \texttt{\_ensure\_schema\_model} 加载解析器模块并缓存该模块。 \texttt{parse\_schema} 用 \texttt{int} 替换 \texttt{int64\_t}，恢复仅关键字参数，并通过 \texttt{\_type\_from\_tg} 将每个解析的类型映射到本地 \texttt{Type} 记录。 \texttt{Type} 存储 \texttt{kind}、\texttt{is\_list}、\texttt{is\_opt}、\texttt{mutability} 和 \texttt{list\_elem\_opt}；助手 \texttt{is\_tensor\_like}、\texttt{is\_mutable\_ref} 和 \texttt{is\_mutable\_tensor\_list} 驱动下游的可变性处理。

\texttt{NativeFunction} 携带 \texttt{schema}、\texttt{base\_name}、\texttt{cpp\_name}、\texttt{overload\_name}、\texttt{args: list[Argument]}、\texttt{returns: list[ReturnDecl]}、\texttt{variants}、\texttt{dispatch}、\texttt{device\_check}、\texttt{structured\_delegate}、\texttt{pydefaults}、\texttt{schema\_kind} (\texttt{functional} / \texttt{inplace} / \texttt{mutable})、\texttt{tags} 和 \texttt{backend\_metadata}。视图元数据（\texttt{is\_view\_op}、\texttt{returns\_view\_of\_input}、\texttt{view\_input\_name}）是通过 \texttt{\_ensure\_view\_metadata} 从 \texttt{tools.autograd.gen\_inplace\_or\_view\_type} 延迟加载的。 \texttt{\_native\_function\_from\_yaml} 将本地 YAML 字段合并到引用函数上，设置 \texttt{cpp\_no\_default\_args}，并在将 \texttt{Composite} 规范化为 \texttt{CompositeExplicitAutograd} 后验证调度映射。 \texttt{NativeFunctionCollection} 包装结果列表，并使用与参考生成器相同的预分组公开 \texttt{grouped\_native\_functions} 和 \texttt{grouped\_view\_functions} ，因此下游生成器看到与规范构建相同的 function/out 和 view/view\_copy/view\_inplace 组。

\subsubsection{\texttt{parse\_schema} 通过模式模型}

\texttt{parse\_schema} 是存储库中 FunctionSchema 的唯一拼写。它委托给本地解析器，同时保留本地不变量。步骤：

\begin{enumerate}[nosep, leftmargin=1.2em]
\item 将解析器的 \texttt{int64\_t} 标准化为 \texttt{int}，而本地 \texttt{Type.kind} 通过 \texttt{\_KIND\_ALIASES} 保留 \texttt{int64\_t}。
\item 调用 \texttt{tgm.FunctionSchema.parse} 并断言往返保真度（： \texttt{str(ts) != schema} 引发）。
\item 从 \texttt{ts.arguments.flat\_kwarg\_only} 恢复仅关键字状态，因为 \texttt{ts.arguments.all} 展平了三个存储桶。
\item 按架构顺序将 \texttt{TensorOptionsArguments} 包装器 (\texttt{expand}) 展开为四个组成部分 \texttt{dtype}/\texttt{layout}/\texttt{device}/\texttt{pin\_memory} 参数。
\item 通过 \texttt{\_type\_from\_tg} 映射每个 \texttt{tgm} 类型，该 \texttt{\_type\_from\_tg} 剥离 \texttt{OptionalType} / \texttt{ListType}，检测 \texttt{list\_elem\_opt}，并应用 \texttt{\_KIND\_ALIASES}（\texttt{SymInt} $\rightarrow$ \texttt{int64\_t} 反转）。
\item 实现可变性：当 \texttt{tgm} 注释被写入（ \texttt{a.annotation.is\_write}）且类型为 \texttt{Tensor} 时，用 \texttt{mutability="a"} 重新创建 \texttt{Type} ，因此 \texttt{is\_mutable\_ref} 和 \texttt{is\_mutable\_tensor\_list} 为 true。
\item 当架构已包含 \texttt{input} 时，对带有数字后缀的面向 Python 的重命名 \texttt{self} $\rightarrow$ \texttt{input} 进行重复数据删除。
\end{enumerate}

\begin{table}[H]
\centering
\caption{\texttt{Type} 字段由 \texttt{\_type\_from\_tg} 生成。}
\label{tab:typefields}
\small
\begin{tabular}{lll}
\toprule
场地 & Type & 意义 \\
\midrule
\texttt{kind} & \texttt{str} & \texttt{Tensor}, \texttt{Scalar}, \texttt{DType}, \texttt{Device}, \texttt{int64\_t}, \texttt{double}, \texttt{bool} \\
\texttt{is\_list} & \texttt{bool} & \texttt{true} 代表 \texttt{Tensor[]}、\texttt{int[]}、\texttt{bool[3]} \\
\texttt{is\_opt} & \texttt{bool} & \texttt{true} 代表 \texttt{Tensor?}、\texttt{int?} \\
\texttt{mutability} & \texttt{str | None} & \texttt{"a"} 表示 \texttt{Tensor(a!)}、\texttt{"a!"} 携带参数 \\
\texttt{list\_elem\_opt} & \texttt{bool} & \texttt{true} 代表 \texttt{Tensor?[]} \\
\texttt{list\_size} & \texttt{int | None} & 固定变量的 \texttt{[N]} 或 \texttt{None} \\
\bottomrule
\end{tabular}
\end{table}

\begin{lstlisting}[caption={\texttt{\_type\_from\_tg} and return-kind classification.}, label=lst:parse_schema, style=pystyle]
def _type_from_tg(t) -> Type:          # model.py:203
    is_opt = isinstance(t, tgm.OptionalType)
    inner = t.elem if is_opt else t
    is_list = isinstance(inner, tgm.ListType)
    elem = inner.elem if is_list else inner
    list_elem_opt = isinstance(elem, tgm.OptionalType)
    # ... kind via _KIND_ALIASES, symint/symbool/symfloat ...
    return make_type(kind, is_list, is_opt, None, symint, ...)

# model.py:313 cpp_return_kind
# 0 ret or void -> "void"; >1 ret -> "tuple"
# 1 ret list    -> "list";  1 ret Tensor(a!) -> "mut_ref" else "value"
\end{lstlisting}

\subsection{编排器：（253行）}

定义协调器。 \texttt{CodegenContext} 保存 \texttt{funcs}、\texttt{native\_by\_opname}、\texttt{derivatives}、\texttt{autocast\_ops}、\texttt{autograd\_ops}（将调度程序名称映射到后向节点名称，\texttt{relu\_} 默认为 \texttt{ReluBackward}）、\texttt{out\_dir}、\texttt{pkg\_out}、\texttt{backend\_indices}、\texttt{reference\_functions} 和 \texttt{written}。 \texttt{register\_generator} 填充全局 \texttt{\_GENERATORS} 字典。每个生成器都是一个普通函数，它接收上下文并调用 \texttt{ctx.write} 或 \texttt{ctx.write\_pkg}。 \texttt{DEFAULT\_TARGETS} 列出了默认的十一个：

\begin{table}[H]
\centering
\caption{生成器和工件。每个都通过 \texttt{model.py} 模式投影读取 YAML。行数为 \texttt{*.py}。}
\label{tab:codegen_generators}
\small
\begin{tabular}{llll}
\toprule
目标 & 模块 & 输出 & 核心入门 \\
\midrule
张量方法 &  & \texttt{.cpp} & \texttt{generate\_header}, \texttt{generate\_cpp} \\
重新调度 &  & \texttt{TensorRedispatchGenerated.h} & \texttt{generate\_redispatch\_header} \\
Autograd节点 &  & \texttt{AutogradNodesGenerated.h} & \texttt{generate\_autograd\_nodes} \\
TPX操作 &  & \texttt{.cpp} & \texttt{cpp} \\
自动微分注册 &  & \texttt{TPXAutogradRegistration.cpp} & \texttt{generate\_autograd\_registration} \\
绑定 &  & \texttt{TensorBindingsGenerated.h} & \texttt{generate\_bindings} \\
自动转换 &  & \texttt{AutocastGenerated.cpp} & \texttt{generate\_autocast\_registration} \\
PythonCAPI &  & \texttt{TensorCPythonGenerated.h} & \texttt{\_gen\_python\_capi} \\
结构化 &  & \texttt{StructuredGenerated.h} & \texttt{\_gen\_structured} \\
就地或查看 &  & \texttt{InplaceOrViewGenerated.h} & \texttt{\_gen\_inplace\_or\_view} \\
Python函数式 &  & \texttt{functional.py} & \texttt{generate\_functional\_py} \\
Pyi存根 &  & \texttt{*.pyi} & \texttt{generate\_pyi} \\
\bottomrule
\end{tabular}
\end{table}

CLI 接受 \texttt{--yaml}、\texttt{--derivatives}、\texttt{--out\_dir}、\texttt{--pkg\_out}、\texttt{--pyi\_template}、\texttt{--pyi\_out} 和 \texttt{--targets}（以逗号分隔，默认为 \texttt{DEFAULT\_TARGETS}）。验证未知目标，通过 \texttt{parse\_native\_yaml} 解析 YAML，通过 \texttt{gen\_autograd.load\_derivatives} 加载衍生品，查询 \texttt{autocast\_registered\_ops}，构建上下文，并分派 \texttt{run\_gen} 按顺序迭代目标。因为每个生成器共享一个 \texttt{NativeFunctionCollection} 和一个 \texttt{derivatives} 字典，所以参数名称、默认值和调度键不能不同。

\subsubsection{ \texttt{CodegenContext}}

\texttt{CodegenContext} 是一个数据类，它通过所有十二个生成器进行单个解析。领域：

\begin{table}[H]
\centering
\caption{\texttt{CodegenContext} 字段。}
\label{tab:codegencontext}
\small
\begin{tabular}{lll}
\toprule
场地 & Type & Role \\
\midrule
\texttt{funcs} & \texttt{NativeFunctionCollection} & 有序列表 + \texttt{grouped\_native\_functions} 视图 \\
\texttt{native\_by\_opname} & \texttt{dict[func\_name $\rightarrow$ NativeFunction]} & 按调度程序名称键入的派生查找 \\
\texttt{derivatives} & \texttt{dict[op $\rightarrow$ OpDerivatives]} & 187 个向后公式 + \texttt{relu\_} 默认值 \\
\texttt{autocast\_ops} & \texttt{set[str]} & 使用 \texttt{CUDA} 内核的操作 \\
\texttt{autograd\_ops} & \texttt{dict[str,str]} & \texttt{func\_name $\rightarrow$ node\_name} \\
\texttt{out\_dir} & \texttt{str} & \texttt{GEN\_INCLUDE\_DIR} 代表 \texttt{ctx.write} \\
\texttt{pkg\_out} & \texttt{str | None} & \texttt{TP\_PACKAGE\_DIR} 代表 \texttt{ctx.write\_pkg} \\
\texttt{backend\_indices} & \texttt{dict} & 用于结构化检查的每键调度索引 \\
\texttt{reference\_functions} & \texttt{tuple} & 用于分组的解析器记录 \\
\texttt{written} & \texttt{dict[str,str]} & 文件名 $\rightarrow$ 新鲜度测试路径 \\
\bottomrule
\end{tabular}
\end{table}

两个属性组函数：\texttt{grouped\_native\_functions} 通过 \texttt{NativeFunctionsGroup.from\_dict} 对函数/输出变体进行分组； \texttt{grouped\_view\_functions} 对别名组 (\texttt{view} / \texttt{view\_copy} / \texttt{view\_inplace}) 执行相同的操作。两个方法抽象I/O：\texttt{write}连接\texttt{out\_dir}，确保父目录，写入并记录\texttt{written}； \texttt{write\_pkg} 在 \texttt{pkg\_out} 下执行相同的操作并断言它已设置。 \texttt{backend\_index} 通过 \texttt{name} 字符串相等性来解析调度键。

\subsubsection{ \texttt{DEFAULT\_TARGETS}}

\texttt{DEFAULT\_TARGETS} 枚举在每个 \texttt{cmake --build} 上运行的十一个目标：

\begin{lstlisting}[caption={default target list.}, label=lst:defaulttargets, style=pystyle]
DEFAULT_TARGETS = ["TensorMethods", "Redispatch", "AutogradNodes", "TPXOps",
                   "AutogradRegistration", "Bindings", "Autocast",
                   "InplaceOrView", "Structured", "PythonFunctional",
                   "PythonCAPI"]
\end{lstlisting}

\texttt{PyiStub} 是第十二个注册的生成器（\texttt{\_gen\_pyi}），但被排除在 \texttt{DEFAULT\_TARGETS} 之外，因为它通过 \texttt{ctx.pyi\_out} 写入源码包目录并需要 \texttt{--pyi\_template} / \texttt{--pyi\_out}。 CLI 默认 \texttt{--targets} 到 \texttt{",".join(DEFAULT\_TARGETS)} 并使用 \texttt{SystemExit(unknown generators)} 验证未知名称。顺序是承载的： \texttt{TensorMethods} before \texttt{Redispatch} before \texttt{TPXOps} 确保在发出 autograd 包装器时存在重新调度标头； \texttt{run\_gen} 按列出的顺序迭代，并通过由 \texttt{register\_generator} 填充的全局 \texttt{\_GENERATORS} 字典调用每个函数。每个生成器都会收到相同的 \texttt{CodegenContext}，因此 YAML 中的重命名会自动传播到所有工件，而无需按目标重新解析。

\subsection{张量曲面：（718 线）}

每个操作员都获得两个表面：一个通过调度程序解析内核的后端重新调度漏斗，以及一个公共 \texttt{Tensor} 成员，该成员在进入重新调度之前执行检查和转换调度。两者均从 \texttt{gen\_api.py} 发出。

\subsubsection{超载计划：\texttt{plan\_members}}

C++ 不能单独重载返回类型、默认值或常量。 \texttt{\_overload\_key} 在剥离 \texttt{requires\_grad} 后将变体折叠为 \texttt{(cpp\_name, tuple(cpp\_arg\_type))}，对于 \texttt{method}，则折叠为 \texttt{self}。当两个成员共享参数类型的公共前缀并且超过该前缀的每个参数都是双方默认的时，\texttt{\_prefix\_ambiguous} 认为两个成员发生冲突；仅提供前缀的调用对于两者都是可行的，并且重载解析将是不明确的。 \texttt{\_HANDWRITTEN\_TENSOR\_METHODS} 列出了（\texttt{dim}、\texttt{numel}、\texttt{detach}、\texttt{as\_strided}、\texttt{to} 等）拥有的成员；手写的无参数成员加上生成的带默认值的单参数成员会产生相同的歧义，因此手写的所有者获胜。 \texttt{plan\_members} 按 YAML 顺序迭代模式，记录每个 \texttt{cpp\_name} 发出的参数列表，并在新变体发生冲突或手写时将 \texttt{plan[(func\_name, variant)]} 标记为 false。 \texttt{generate\_header} 和 \texttt{generate\_cpp} 都会查阅该地图，因此被抑制的成员仍然保留其重新调度帮助程序和注册条目。

\begin{table}[H]
\centering
\caption{\texttt{plan\_members} 决策表。}
\label{tab:planmembers}
\small
\begin{tabular}{lll}
\toprule
健康）状况 & Test & 影响 \\
\midrule
手写业主 & \texttt{cpp\_name in \_HANDWRITTEN\_TENSOR\_METHODS} & \texttt{plan[func,variant]=False} \\
前缀歧义 & \texttt{\_prefix\_ambiguous(prev, new)} & 后来的变种被抑制 \\
精确类型冲突 & \texttt{\_overload\_key} 已经在 \texttt{seen\_decl} 中 & 重复数据删除，不排放 \\
否则 & YAML顺序，先胜 & \texttt{plan[func,variant]=True} \\
\bottomrule
\end{tabular}
\end{table}

\subsubsection{标头发射：\texttt{generate\_header}}

\texttt{generate\_header} 将 \texttt{TensorGenerated.h} 写在 \texttt{\#pragma once} 下，并带有 \texttt{\#include <tuple>}。它重播 \texttt{plan\_members}，对 \texttt{\_overload\_key} 进行重复数据删除，并为 \texttt{function} 变体（工厂表面 \texttt{Tensor::zeros}）和实例成员发出 \texttt{static}，否则，当 \texttt{self} 不是 \texttt{a!} 时，通过 \texttt{\_is\_const\_method} 附加 \texttt{const}。签名由 \texttt{method\_signature} 使用 \texttt{api\_types.cpp\_return\_type} 和 \texttt{cpp\_arg\_type} 渲染，默认值来自 \texttt{cpp\_default}，除非 \texttt{cpp\_no\_default\_args} 抑制它们或操作是 \texttt{.out}。结果是一个标头被设计为包含在 \texttt{class Tensor} (\S\ref{sec:codegen_class_mixing}) 内，而不是包含在命名空间范围内。

\subsubsection{设备和密钥解析：\texttt{\_device\_source} 和 \texttt{\_dispatch\_key\_expr}}

\texttt{\_device\_source} 对于一种架构和变体返回 \texttt{(dispatch\_key\_source, target\_device\_expr)}。如果存在 \texttt{device} 参数，则直接使用它，对于 \texttt{\_like} 工厂，可选设备回退到 \texttt{Device(DeviceType::CPU)} 或 \texttt{self.device} ；对于 \texttt{method} 变体，\texttt{device} 表示 \texttt{self.device}，稍后在自由函数重新调度中重写为 \texttt{self.device}。否则第一个非可选 \texttt{Tensor}，然后通过 \texttt{deviceForTensorArg} 第一个 \texttt{Tensor[]}，否则为 CPU。 \texttt{\_dispatch\_key\_expr} 选择最高的活动变换键：\texttt{copy\_} 使用其第一个参数，\texttt{\_RANDOM\_TRANSFORM\_OPS} 中的随机工厂使用 \texttt{transform::dispatch\_key\_for\_random(computeDispatchKey(device))}，并且一般情况会在每个类似张量的参数上折叠 \texttt{dispatchKeyForTensorArgs(self, a, b)} （方法接收器在重新调度外部变为 \texttt{*this}，在内部变为 \texttt{self}）。

\subsubsection{重新调度漏斗：\texttt{\_emit\_redispatch} 和 \texttt{OpRecord}}

\texttt{\_emit\_redispatch} 在 \texttt{namespace detail} 下为每个模式发出一个自由函数 \texttt{detail::redispatch\_<op>\_<variant>}：

\begin{lstlisting}[caption={skeleton of \texttt{\_emit\_redispatch} (simplified).}, label=lst:emitredispatch, style=cppstyle]
namespace detail {
TENSORPLAY_API Tensor redispatch_add_Tensor_method(const Tensor& self,
                                                   const Tensor& other, Scalar alpha) {
    tensorplay::prof::OpRecord __tp_prof_rec("add.Tensor");
    if (tensorplay::prof::g_capture_shapes.load(std::memory_order_acquire)) {
        // capture shapes/dtypes of Tensor-like args into __tp_prof_shapes
    }
#ifdef USE_CUDA
    tensorplay::prof::GpuTimerPair __tp_gpu(__tp_prof_rec);
#endif
#ifdef USE_CUDA
    cuda::OptionalCUDAGuard device_guard(self.device());
#endif
    static const OperatorHandle op_handle =
        Dispatcher::singleton().findHandle("add.Tensor");
    DispatchKey dispatch_key = toBackendKey(dispatchKeyForTensorArgs(self));
#ifdef USE_CUDA
    __tp_gpu.arm(self.device());
#endif
    auto&& __tp_result = DispatchStub<Tensor,const Tensor&,const Tensor&,Scalar>
        ::call(op_handle, dispatch_key, self, other, alpha);
#ifdef USE_CUDA
    __tp_gpu.close();
#endif
    // output allocation volume and version bump handled below
    return std::move(__tp_result);
}
} // namespace detail
\end{lstlisting}

具体来说，发射器写入：由 \texttt{func\_name} 键控的 \texttt{OpRecord}，由 \texttt{g\_capture\_shapes} 保护的形状捕获分支，枚举 \texttt{Tensor}、\texttt{Tensor[]}、\texttt{Tensor?} 和 \texttt{Tensor?[]}，一个 \texttt{GpuTimerPair} 池支持的 \texttt{cudaEvent} 对，当 \texttt{device\_guard} 为 true 时的 \texttt{OptionalCUDAGuard}，通过 \texttt{Dispatcher::singleton.findHandle(func\_name)} 解析一次的 \texttt{static const OperatorHandle}，由 \texttt{\_dispatch\_key\_expr} 计算出的 \texttt{DispatchKey} 和用于非随机操作的 \texttt{toBackendKey}，在分派前有 \texttt{\_\_tp\_gpu.arm(device)} ，在分派后有 \texttt{\_\_tp\_gpu.close} ，通过句柄键入的 \texttt{DispatchStub<Ret,Args...>::call} ，以及每个可变参数的版本碰撞：由 \texttt{!InferenceMode::is\_enabled} 保护的 \texttt{unsafeGetTensorImpl->bump\_version} ，在 \texttt{Tensor(a!)} 数组上有列表可变循环。主体是分析器的单个每操作检测点。

\begin{table}[H]
\centering
\caption{\texttt{\_emit\_redispatch} 发射顺序。}
\label{tab:emitredispatch}
\small
\begin{tabular}{lll}
\toprule
Step & 线路 & Code \\
\midrule
1 个操作记录 &  & \texttt{OpRecord \_\_tp\_prof\_rec("func\_name")} \\
2 形状捕捉 & --\texttt{389} & \texttt{g\_capture\_shapes} $\rightarrow$ \texttt{set\_io\_meta(shapes,dtypes)} \\
3 Gpu定时器对 &  & \texttt{GpuTimerPair \_\_tp\_gpu(\_\_tp\_prof\_rec)} \\
4 设备防护 &  & \texttt{OptionalCUDAGuard(device)} \\
5 手柄 &  & \texttt{static OperatorHandle = findHandle(func\_name)} \\
6 调度密钥 &  & \texttt{toBackendKey(dispatchKeyForTensorArgs(\dots))} \\
7 派遣 &  & \texttt{DispatchStub<Ret,Args>::call(handle,key,args)} \\
8 版本碰撞 &  & \texttt{bump\_version} 每 \texttt{Tensor(a!)} \\
\bottomrule
\end{tabular}
\end{table}

\subsubsection{热心会员：\texttt{generate\_cpp} 和设备卫士}

\texttt{generate\_cpp} 运行两次。 Pass~1 发出在 \texttt{redispatch\_key} 上进行重复数据删除的重新调度渠道。 Pass~2 发出公共 \texttt{Tensor} 成员，在 \texttt{\_overload\_key} 上进行重复数据删除并通过 \texttt{plan\_members} 进行过滤。每个成员都是合格的 \texttt{Tensor::} 并以特殊情况的快速路径打开：\texttt{copy\_} 验证 \texttt{shape} 相等性，\texttt{contiguous} 解析 \texttt{MemoryFormat::Preserve} 并在已经连续时返回 \texttt{*this} 以避免用 \texttt{ContiguousBackward} 节点标记叶别名。 \texttt{\_emit\_device\_checks} 如下：标量 \texttt{Tensor} 和 \texttt{Tensor?} 发出 \texttt{if (a.device != target) TP\_THROW(DeviceMismatchError,\dots)}； \texttt{Tensor[]} 循环元素并报告有问题的参数名称。对于目标设备为目标的 \texttt{device\_check == NoCheck} 和 \texttt{copy\_} ，将跳过检查。

成员内部键的顺序是承重的。设备检查后，首先探测偏移量~9处的vmap：计算\texttt{\_\_tp\_key = dispatchKeyForTensorArgs(...)}，测试\texttt{is\_vmap\_key}，如果没有注册批处理规则，则通过带有\texttt{NotImplementedError}的vmap句柄直接\texttt{DispatchStub::call}。 Autocast 是第二个：计算 \texttt{\_\_ac\_key = toAutocastKey(computeDispatchKey(dev\_src))}，在探测 \texttt{getKernel} 之前检查 \texttt{autocast::is\_enabled}，如果存在则调度。 Autograd 是第三个：\texttt{autograd\_ops} 将 \texttt{func\_name} 映射到后向节点名称，并且当 \texttt{GradMode::is\_enabled} 和内核注册时，成员通过 \texttt{toAutogradKey} 调度。如果没有任何转换键触发，则尾部调用 \texttt{detail::redispatch\_<op>\_<variant>}。

\subsubsection{重新调度标头：\texttt{generate\_redispatch\_header}}

\texttt{generate\_redispatch\_header} 将 \texttt{TensorRedispatchGenerated.h} 与 \texttt{\#include "Tensor.h"} 和 \texttt{\#include "Macros.h"} 写入 \texttt{namespace tensorplay::detail} 内，并声明每个 \texttt{TENSORPLAY\_API Ret redispatch\_<op>\_<variant>(stub\_args)} 一次。该标头由 eager 编译单元和 autograd 包装器包含，以便 autograd 层可以调用相同的分派内核，而无需重新解析句柄。

\subsection{Autograd 曲面：（624 线）}

\texttt{gen\_tpx.py} 产生位于 \texttt{tpx} 中的 \texttt{AutogradCPU}/\texttt{AutogradCUDA} 内核。每个算子都通过自动转换、突变防护、后端和后向节点构造来公开 \texttt{tensorplay::tpx::ops::<name>} 路由。

\subsubsection{内联谓词：\texttt{\_inline\_forward\_lines} vs.\ 非内联}

\texttt{\_inline\_forward\_lines} 决定包装器是否不进行急切路径工作并且可以作为内联标头定义发出。当操作符没有向后（\texttt{\_has\_autograd}，包括手动关闭的 \texttt{relu\_}）、没有自动转换规则、没有 \texttt{device} 或 \texttt{dtype} 的工厂全局默认解析（当包装器必须填充 \texttt{DType} 时）、没有视图版本计数器共享以及没有 \texttt{requires\_grad} 管道时，它返回一个主体。否则，它返回 \texttt{None} 并强制在 \texttt{TPXOpsGenerated.cpp} 中出现一个外线主体。核心调用表达式 \texttt{\_core\_call\_expr} 选择 \texttt{::tensorplay::detail::redispatch\_<op>\_<variant>}，或者对于 \texttt{manual\_kernel\_registration}，选择底层 \texttt{Tensor::} 成员。

\begin{table}[H]
\centering
\caption{内联谓词与非内联谓词。}
\label{tab:tpxinline}
\small
\begin{tabular}{lll}
\toprule
健康）状况 & Test & 结果 \\
\midrule
有落后 & \texttt{\_has\_autograd(f,derivatives)} & 非内联 \\
多输出视图 & \texttt{list return + returns\_view\_of\_input + dim} & 非内联 \\
自动转换规则 & \texttt{func\_name in autocast\_ops} & 非内联 \\
出厂设备默认值 & \texttt{device: optional<Device>} 无张量 & 非内联 \\
工厂默认数据类型 & \texttt{dtype in \_DTYPE\_GLOBAL\_DEFAULT\_OPS} & 非内联 \\
查看分享 & \texttt{returns\_view\_of\_input + value} & 非内联 \\
需要等级管道 & \texttt{requires\_grad arg + value} & 非内联 \\
否则 & 普通转发器 & 排队 \\
\bottomrule
\end{tabular}
\end{table}

\begin{lstlisting}[caption={Inline vs.\ non-inline emission.}, label=lst:tpxinline, style=cppstyle]
// gen_tpx.py:137 inline trivial forwarder in TPXOpsGenerated.h
inline Tensor add(const Tensor& self, const Tensor& other, Scalar alpha = 1) {
    return ::tensorplay::detail::redispatch_add_Tensor_function(self, other, alpha);
}
// gen_tpx.py:140 non-inline declaration
TENSORPLAY_API Tensor mm(const Tensor& self, const Tensor& mat2);
// body lives in TPXOpsGenerated.cpp with 5 blocks
\end{lstlisting}

\subsubsection{标头拆分：\texttt{generate\_tpx\_ops\_h}}

\texttt{generate\_tpx\_ops\_h} 对 \texttt{\_dedup\_key}（C++ 名称加上参数类型列表）进行重复数据删除，计算 \texttt{cpp\_return\_type} 和 \texttt{\_signature\_args}，对于 \texttt{inline} 操作直接发出 \texttt{inline Ret name(args) \{ return detail::redispatch(...); \}}；其余的它发出 \texttt{TENSORPLAY\_API Ret name(args);} 作为前向声明。这种分割将热的普通转发器保留在标头中以进行内联，同时将图形构建包装器保留在标头之外。

\subsubsection{包装体：\texttt{generate\_tpx\_ops\_cpp}}

每个外线包装器都遵循相同的五个块。 (1) 工厂全局默认分辨率从 \texttt{globalContext.defaultDevice} 填充 \texttt{device}，对于 \texttt{rand}、\texttt{randn}、\texttt{empty}、\texttt{zeros}、\texttt{ones}、\texttt{dtype}，当调用者未设置它们时，从 \texttt{globalContext.defaultDType} 填充； \texttt{arange}/\texttt{linspace} 被排除，因为它们从标量输入推断数据类型。 (2) 自动转换：\texttt{\_emit\_autocast\_block} 从第一个张量参数重新导出设备源，并在 \texttt{autocast::is\_enabled} 时通过 \texttt{toAutocastKey} 进行调度。 (3)梯度簿记：\texttt{\_emit\_requires\_grad\_detection}从\texttt{GradMode}、\texttt{InferenceMode}以及任何需要梯度的张量中设置\texttt{bool requires\_grad}； \texttt{\_emit\_leaf\_checks} 禁止叶或多输出视图的就地突变；对于派生使用 \texttt{self} 的就地操作，突变前克隆 \texttt{\_\_tp\_original\_self} 在 \texttt{GradMode} 禁用下捕获。 (4) 核心调用按 \texttt{cpp\_return\_kind} (\texttt{value}、\texttt{tuple}、\texttt{list}、\texttt{mut\_ref}、\texttt{void}) 存储 \texttt{core\_result} / \texttt{result}。对于视图操作，包装器共享版本计数器，标记 \texttt{is\_view}，并通过重播 lambda \texttt{\_view\_replay\_lambda} 使用 \texttt{CreationMeta::DEFAULT} 或 \texttt{MULTI\_OUTPUT\_NODE} 安装视图元数据。 (5) 后向节点：\texttt{\_node\_ctor\_args} 从 \texttt{derivatives} 公式中排序构造函数参数，\texttt{\_emit\_edges} 收集 \texttt{Edge}（包括列表类型重载），并且结果按每个返回类型用 \texttt{set\_grad\_fn} / \texttt{set\_requires\_grad} 包装。 \texttt{generate\_autograd\_registration} 然后通过 \texttt{Dispatcher::registerKernel} 为键 \texttt{AutogradCPU}、\texttt{AutogradCUDA}、\texttt{AutogradVulkan} 注册每个包装器一次。

\subsection{剩余发电机}

\texttt{generate\_autograd\_nodes} 通过表达式 AST（分词器加优先级攀爬解析器）而不是正则表达式进行解析，遍历 AST 一次以计算保存的变量集，并写入 \texttt{AutogradNodesGenerated.h}，每个向后节点使用一个结构来保存保存的张量和 \texttt{apply} 主体。 \texttt{generate\_bindings} 写入 \texttt{TensorBindingsGenerated.h} 通过 \texttt{pybind11} 暴露 \texttt{bind\_generated\_tensor\_methods} 和 \texttt{bind\_generated\_op\_functions}，跳过 FASTCALL 层已声明的任何操作。 \texttt{generate\_autocast\_registration} 写入 \texttt{AutocastGenerated.cpp} 注册自动转换内核。并为结构化内核和别名/就地元数据编写单头填充程序；两者都包含在构建的早期，因此调度程序在 Python 扩展链接之前看到它们的内核。

\subsection{构建分期：和}

构建文件定义了暂存合同。 \texttt{GEN\_BASE\_DIR} 默认为 \texttt{generated}，\texttt{GEN\_INCLUDE\_DIR} 默认为 \texttt{ops}。十个变量枚举阶段输出：\texttt{GEN\_HEADER} (\texttt{TensorGenerated.h})、\texttt{GEN\_CPP}、\texttt{GEN\_BINDINGS\_HEADER}、\texttt{GEN\_CPYTHON\_HEADER}、\texttt{GEN\_AUTOGRAD\_NODES\_H}、\texttt{GEN\_TPX\_OPS\_H}、\texttt{GEN\_TPX\_OPS\_CPP}、\texttt{GEN\_TPX\_AUTOGRAD\_REG\_CPP}、\texttt{GEN\_AUTOCAST\_CPP} 和 \texttt{GEN\_TPX\_REDISPATCH\_H}。同一个文件声明一个 \texttt{add\_custom\_command}，其 \texttt{OUTPUT} 是上述所有八个加上绑定头，其\texttt{COMMAND} 是 \texttt{derivatives.yaml --out\_dir GEN\_INCLUDE\_DIR --pkg\_out tensorplay}，\texttt{PYTHONPATH} 是存储库的根目录，其 \texttt{DEPENDS} 列出了两个 YAML 文件、所有 \texttt{*.py} 模块和 \texttt{.cpp}，因此桥接更改会触发重新生成。它标记每个输出 \texttt{GENERATED TRUE} 并创建 \texttt{p10}、\texttt{tpx} 和 \texttt{\_C} 所依赖的虚假目标 \texttt{generate\_code}；生成文件夹中的包含目录使生成的标头可解析。

该规则还暂存包输出。将 \texttt{TP\_PACKAGE\_DIR} 设置为 \texttt{tensorplay} 并且生成器通过 \texttt{gen\_pyi} 将 \texttt{functional.py} （见下文）和 \texttt{*.pyi} 写入源包目录，因此普通的 \texttt{pip install} 无需单独的生成步骤即可传送它们。 \texttt{p10} 选择退出 \texttt{UNITY\_BUILD} （内核翻译单元必须保持独立以进行静态注册），而 \texttt{\_C} 选择加入以摊销 9\,300 行 \texttt{TensorCPythonGenerated.h} 和绑定单元。

\begin{lstlisting}[caption={generation rule (abridged).}, label=lst:cmake, style=cppstyle]
add_custom_command(
  OUTPUT ${GEN_HEADER} ${GEN_CPP} ${GEN_BINDINGS_HEADER}
         ${GEN_CPYTHON_HEADER} ${GEN_AUTOGRAD_NODES_H}
         ${GEN_TPX_OPS_H} ${GEN_TPX_OPS_CPP}
         ${GEN_TPX_AUTOGRAD_REG_CPP} ${GEN_AUTOCAST_CPP}
         ${GEN_TPX_REDISPATCH_H}
  COMMAND ${CMAKE_COMMAND} -E env PYTHONPATH=${CMAKE_CURRENT_SOURCE_DIR}
          ${Python_EXECUTABLE} -m tools.codegen.main
            --yaml ${CMAKE_CURRENT_SOURCE_DIR}/config/native_functions.yaml
            --derivatives ${CMAKE_CURRENT_SOURCE_DIR}/config/derivatives.yaml
            --out_dir ${GEN_INCLUDE_DIR}
            --pkg_out ${CMAKE_CURRENT_SOURCE_DIR}/tensorplay
  DEPENDS config/native_functions.yaml config/derivatives.yaml
          tools/codegen/main.py tools/codegen/model.py
          tools/codegen/gen_api.py tools/codegen/gen_tpx.py
          tools/codegen/gen_python.py tools/codegen/gen_pyi.py
           tools/codegen/*.py
)
set_source_files_properties(${GEN_HEADER} PROPERTIES GENERATED TRUE)
\end{lstlisting}

\subsubsection{解析器依赖分段}

模式解析器通过 \texttt{PYTHONPATH} 在树内使用。构建显式依赖集：

\begin{lstlisting}[caption={parser collection.}, label=lst:cmakeparser, style=cppstyle]
file(GLOB_RECURSE tp_codegen_parser CONFIGURE_DEPENDS
    "${CMAKE_CURRENT_SOURCE_DIR}/tools/codegen/*.py")
\end{lstlisting}

每个匹配的文件都会添加到单个 \texttt{add\_custom\_command} 的 \texttt{DEPENDS} 中，因此触摸解析器或生成器模块会强制在下一个 \texttt{ninja -C build} 链接之前重新生成。 \texttt{COMMAND} 设置 \texttt{PYTHONPATH=\$\{CMAKE\_CURRENT\_SOURCE\_DIR\}}，因此本地架构模块无需安装步骤即可解析。生成的标头被标记为 \texttt{GENERATED TRUE}，在 \texttt{clean} 上删除，并被 git via 忽略。领带订购的虚假目标：

\begin{table}[H]
\centering
\caption{CMake 构建暂存依赖项 (--\texttt{923})。}
\label{tab:cmakedep}
\small
\begin{tabular}{lll}
\toprule
变量/目标 & 价值 & 消费者 \\
\midrule
\texttt{GEN\_BASE\_DIR} & \texttt{generated} & \texttt{include\_directories(GEN\_BASE\_DIR)} \\
\texttt{GEN\_INCLUDE\_DIR} & \texttt{ops} & 10 个标头 + 2 个来源 \\
\texttt{TP\_PACKAGE\_DIR} & \texttt{tensorplay} & \texttt{functional.py}, \texttt{*.pyi} \\
\texttt{tp\_codegen\_parser} & \texttt{*.py} & \texttt{DEPENDS} 触发再生 \\
\texttt{generate\_code} & 假 & \texttt{p10}、\texttt{tpx}、\texttt{\_C} 取决于 \\
\texttt{.cpp} & 桥 & \texttt{DEPENDS} 所以桥梁变化重新生成 \\
\bottomrule
\end{tabular}
\end{table}

\subsection{类混合：和两个标题}
\label{sec:codegen_class_mixing}

将类 \texttt{Tensor} 定义为 \texttt{shared\_ptr<TensorImpl>} 句柄。它包括生成的成员声明：

\begin{lstlisting}[caption={mixing site.}, label=lst:mixing, style=cppstyle]
class P10_API Tensor {
  // ... dim(), numel(), device(), is_contiguous(), manual members ...
#include "tensorplay/ops/TensorGenerated.h"
};
\end{lstlisting}

由于包含发生在类定义内部，因此枚举的 \texttt{static Tensor zeros(...)} 成员成为 \texttt{Tensor} 的静态成员，而 \texttt{Tensor add(const Tensor\& other) const} 成为具有正确 \texttt{const} 限定符的实例方法。生成的标头仅提供声明；匹配定义位于 \texttt{TensorGenerated.cpp} 中，其中包括 \texttt{Tensor.h} 和 \texttt{TensorRedispatchGenerated.h}，并作为 \texttt{p10} 的一部分进行编译。第二个生成的标头位于 \texttt{namespace detail} 中，并由 \texttt{p10} 和 \texttt{tpx} 包含，为 autograd 层提供了一个稳定的名称来调用，而无需重新解析 YAML。

这种分割保证了单一所有权规则：\texttt{\_HANDWRITTEN\_TENSOR\_METHODS} 中列出的手动成员永远不会被重新声明，不明确的重载会在生成时而不是编译时被修剪，并且每个剩余的声明都只有一个定义。

\subsubsection{\texttt{TensorGenerated.h} 混合细节}

混合合约仅包含标头，并依赖于在生成时检查的三个不变量：

\begin{table}[H]
\centering
\caption{\texttt{TensorGenerated.h} 混合不变量。}
\label{tab:tensorh_mixing}
\small
\begin{tabular}{lll}
\toprule
不变式 & 执行 & 引文 \\
\midrule
包括内部班级 & \texttt{TensorGenerated.h"} at & \\
静态工厂 & \texttt{variant == "function"} 时的 \texttt{static} 前缀 & \\
常量方法 & 当 \texttt{self} 不是 \texttt{a!} 时，\texttt{const} 后缀 & \texttt{\_is\_const\_method} \\
无需重新申报 & \texttt{plan\_members} + \texttt{\_HANDWRITTEN\_TENSOR\_METHODS} & , \texttt{64} \\
仅标头声明 & \texttt{\#pragma once} + \texttt{\#include <tuple>} & \\
杂注一次 & 每个工件都以 \texttt{ Generated -- DO NOT EDIT} 开头 & , \texttt{690} \\
\bottomrule
\end{tabular}
\end{table}

\texttt{generate\_header} 发出 \texttt{\#pragma once} 和 \texttt{\#include <tuple>} 一次，然后迭代 YAML 顺序，跳过 \texttt{plan\_members} 为 false 的任何 \texttt{(func\_name, variant)}，对 \texttt{\_overload\_key} 进行重复数据删除，并通过使用 \texttt{api\_types.cpp\_return\_type} 和 \texttt{cpp\_arg\_type} 的 \texttt{method\_signature} 进行渲染。配套的 \texttt{TensorGenerated.cpp} 是独立的：它包括 \texttt{Tensor.h} 和 \texttt{TensorRedispatchGenerated.h} ，并被编译为 \texttt{p10} 的一部分（它选择退出 \texttt{UNITY\_BUILD} 因此静态内核注册不会折叠）。第二个标头 \texttt{TensorRedispatchGenerated.h} 命名空间为 \texttt{tensorplay::detail} ，并为每个 \texttt{redispatch\_key} 声明一次 \texttt{TENSORPLAY\_API Ret redispatch\_<op>\_<variant>(args)} ，从而为 \texttt{tpx} 提供了稳定的调用目标，而无需重新解析。

\subsection{类型映射：（540 行）}

\texttt{api\_types.py} 是 C++ 拼写的单一词汇表。 \texttt{cpp\_arg\_type} 将 \texttt{Type} 映射到 \texttt{gen\_api.py} 中看到的存根类型：\texttt{Tensor(a!)} 变为 \texttt{Tensor\&}，可选列表变为 \texttt{std::optional<std::vector<int64\_t>>}，\texttt{SymInt} 别名到 \texttt{int64\_t} ABI 边界。 \texttt{api\_types.cpp\_return\_type} 通过检查 \texttt{NativeFunction.returns} 和 \texttt{cpp\_return\_kind} 将返回分为 \texttt{void}、\texttt{value}、\texttt{tuple}、\texttt{list} 或 \texttt{mut\_ref}；这会驱动标头返回类型和重新调度存根模板。 \texttt{binding\_default} 和 \texttt{cpp\_default} 渲染模式默认用于 C++ 标头与 pybind11 绑定，而 \texttt{python\_default} 为 Python 表面渲染 \texttt{None}/\texttt{True}/\texttt{False} 。该模块由所有三个 C++ 生成器导入，因此对 \texttt{MemoryFormat} 或 \texttt{Device} 的拼写的更改会自动传播。

\subsection{导数编译：（785行）}

导数不是字符串替换的；它们被编译。 \texttt{gen\_autograd.py} 标记每个公式，构建表达式 AST（节点 \texttt{Num}、\texttt{BoolLit}、\texttt{Var}、\texttt{Call}），并运行优先级攀爬解析器，将比较转换为分派操作（\texttt{\_normalize\_comparisons} 将 \texttt{a > b} 重写为 \texttt{gt(a,b)}）。 \texttt{OpDerivatives} 存储 \texttt{node\_name}、\texttt{members}、\texttt{used\_input\_names} 和每个输出公式。对 AST 的一次遍历会收集保存的变量需求，因此结构定义、其构造函数参数列表以及构造 \texttt{grad\_fn} 的每个调用站点都在构造上一致。无法表示为标量公式的向后列表映射（\texttt{cat}、\texttt{stack}、\texttt{roll}）在 \texttt{MANUAL\_DERIVATIVES} 中声明并作为手写节点实例发出。输出 \texttt{AutogradNodesGenerated.h} 包含 148 个节点结构；每个都在需要时将保存的张量保存为 \texttt{SavedVariable} ，并保存一个 \texttt{apply} ，以正确的广播重播公式。

\subsection{Autocast、结构化和 Alias 垫片}

读取与 \texttt{gen\_api.py} 和 \texttt{gen\_tpx.py} 测试相同的 \texttt{autocast\_registered\_ops} 集，并写入 \texttt{AutocastGenerated.cpp}，通过 \texttt{Dispatcher::registerKernel} 将每个自动转换内核注册为 \texttt{DispatchKey::AutocastCPU}/\texttt{AutocastCUDA}。该文件是 \texttt{GENERATED TRUE} 并链接到 \texttt{p10}，但仅当 \texttt{autocast::is\_enabled} 为 true 时才会查阅其内核；禁用的路径是单线程本地负载。每个人都写一个单头垫片。结构化委托记录哪些操作转发到元内核；别名填充程序公开视图谓词，以便 Python 表面可以报告 \texttt{is\_view} 而无需接触 C++ 对象。两者都故意很小：它们的内容仅是标头，并且贡献零符号，但它们在 \texttt{DEPENDS} 中的存在使别名表与 YAML 保持同步。

\subsection{Python Surface：（796 行）和}

\texttt{generate\_functional\_py} 写入 \texttt{functional.py}。标头导入 \texttt{tensorplay.\_C as \_C} 和 \texttt{tensorplay.graph.capture\_call}；帮助程序 \texttt{\_ensure\_device} 和 \texttt{\_MISSING} 规范化 \texttt{device} 拼写，绑定无法解析和保留哨兵区分的可选参数。对于每个 \texttt{cpp\_name}，生成器收集重载组，跳过关键字名称，跳过 \texttt{linalg\_*}（通过 \texttt{tensorplay.linalg} 公开），并应用一系列特殊情况：\texttt{where}（标量提升）、\texttt{div}/\texttt{divide}（可选 \texttt{rounding\_mode}）、\texttt{\_foreach\_*}（可变参数）、 \texttt{mm}/\texttt{addmm}/\texttt{matmul}/\texttt{linear}（图形捕获放置）、工厂 \texttt{randn}、\texttt{rand}、\texttt{zeros}、\texttt{ones}、\texttt{empty} 以及 \texttt{pin\_memory} 快速路径缓存 \texttt{\_c\_<name>}、\texttt{arange}（从浮点参数推断数据类型）、\texttt{squeeze\_copy}、\texttt{unique\_consecutive}、 \texttt{trapz} 和 \texttt{kaiser\_window}。一般路径处理缩减合并：\texttt{\_reduction\_union\_lines} 将 \texttt{sum} 加 \texttt{sum.dim\_IntList} 折叠成一个 \texttt{def sum(input, dim=None, keepdim=False, *, dtype=None)}，将 \texttt{dim is None} 路由到基础内核，否则路由到暗淡内核。多重载组按位置向前移动，以便 \texttt{METH\_FASTCALL} 组调度员可以按形状挑选候选人；单重载组按关键字转发以提高可读性并应用 Scalar/\texttt{int[]} 加宽。每个包装器都会探测 \texttt{\_capturing}，并且当 \texttt{Tracer} 处于活动状态时，会通过 \texttt{\_capture\_call} 进行转向；否则它会调用 \texttt{\_C.<name>} 且 \texttt{device} 已标准化。

\texttt{gen\_pyi.py} 写入类型存根。 \texttt{generate\_pyi} 按 \texttt{function} vs.\ \texttt{method} vs.\ \texttt{functional} 对提示进行分组，将 \texttt{.out} 变体合并到可选的 \texttt{out: Tensor | None = None} 参数中，将列表默认值呈现为 \texttt{\_size} vs.\ \texttt{\_int}，并通过 \texttt{gather\_docstrs} 收集运行时文档字符串，而无需导入包（\texttt{Mock} 记录 \texttt{\_add\_docstr} 调用）。 \texttt{\_dtype\_block} 从 \texttt{DType.h} 或 \texttt{DType.cpp} 派生 \texttt{DType} 枚举成员，因此存根永远不会偏离运行时枚举。三个输出通过 \texttt{*.pyi.in} 模板中的 \texttt{\_substitute} 渲染：\texttt{\_\_init\_\_.pyi}、\texttt{\_VariableFunctions.pyi} 和 \texttt{functional.pyi}。

\begin{lstlisting}[caption={and a stub signature (abridged).}, label=lst:functionalpyi, style=pystyle]
# tensorplay/functional.py (generated, 950 lines)
def add(input, other, alpha=1):
    if _capturing():
        _captured = _capture_call(add, (input, other, alpha), {})
        if _captured is not None:
            return _captured
    if not isinstance(other, (Tensor, Scalar)):
        other = Scalar(other)
    return _C.add(input, other, alpha=alpha)

def sum(input, dim=None, keepdim=False, *, dtype=None):
    # dim=None union wrapper folding sum + sum.dim_IntList
    if _capturing():
        _captured = _capture_call(sum, (input, dim, keepdim), {'dtype': dtype})
        if _captured is not None:
            return _captured
    if dim is None:
        return _C.sum(input, dtype=dtype)
    return _C.sum(input, dim, keepdim, dtype=dtype)
\end{lstlisting}

\begin{lstlisting}[caption={Stub rendering: \texttt{\_signature} and a generated \texttt{.pyi} line.}, label=lst:pyisig, style=pystyle]
# generated_methods: one @overload per schema behind a *
@overload
def add(self, other: TensorBase, alpha: Number | _complex = 1) -> TensorBase: ...
# generated_functions: out merged as optional keyword
def sum(input: TensorBase, dim: _int | _size | None = None,
        keepdim: _bool = False, *, dtype: _dtype | None = None,
        out: TensorBase | None = None) -> TensorBase: ...
\end{lstlisting}

生成 \texttt{Tensor add(Scalar alpha=1) const;} 的相同 \texttt{NativeFunction} 也会生成上面的两个列表，因此向 YAML 添加参数会在一次 \texttt{python -m tools.codegen.main} 调用中更新 C++ API、Python 包装器和类型检查器。

\subsubsection{Python \texttt{pyi} 生成：}

\texttt{gen\_pyi.py}（586 行）通过块感知替换呈现 \texttt{*.pyi.in} 模板中的三种类型存根。管道：

\begin{table}[H]
\centering
\caption{发电管道。}
\label{tab:pyi}
\small
\begin{tabular}{lll}
\toprule
阶段 & 功能 & 细节 \\
\midrule
团体 & \texttt{\_group\_hints} & 按 \texttt{function} / \texttt{method} / \texttt{functional} 分区 \\
签名 & \texttt{\_signature} & \texttt{Type -> atom} 通过 \texttt{\_PYI\_ATOMS}、\texttt{is\_opt -> | None} \\
输出合并 &  & \texttt{.out} 折叠到 \texttt{out: TensorBase | None = None} \\
超载 &  & 每个模式一个 \texttt{@overload}； \texttt{int[]} 领先列表的 vararg \\
文档字符串 & \texttt{gather\_docstrs} & 模拟 \texttt{\_add\_docstr}，不导入包 \\
数据类型 & \texttt{\_dtype\_block} & \texttt{DType.cpp} \texttt{.value} 或 \texttt{DType.h} 枚举的正则表达式 \\
使成为 & \texttt{\_substitute} & \texttt{\$\{name\}} 块缩进替换 \\
输出 & \texttt{\_PYI\_OUTPUTS} & \texttt{\_\_init\_\_.pyi}, \texttt{\_VariableFunctions.pyi}, \texttt{functional.pyi} \\
\bottomrule
\end{tabular}
\end{table}

\texttt{\_signature} 使用三个规则来格式化 \texttt{def name(self, arg: type = default) -> ret:...}： (1) 方法变体前置 \texttt{self}，函数变体使用公共拼写 \texttt{input} 作为杂散 \texttt{self} 槽； (2) 列表种类加宽：\texttt{int[] -> \_int | \_size} 当默认为标量时，\texttt{Tensor[] -> tuple[TensorBase,...] | list[TensorBase]}； (3) 仅当非输出基址和 \texttt{.out} 同级共享相同的位置参数 (\texttt{\_args\_equal}) 时，\texttt{out} 才会作为可选关键字合并。 \texttt{pyi\_type} 通过 \texttt{\_append\_optional} 追加 \texttt{| None}，这会折叠 \texttt{Device} 的双重子句。 \texttt{\_substitute} 替换 \texttt{\$\{name\}} 孔：当孔开始一行时，其内容通过 \texttt{textwrap.indent} 缩进到孔的深度，因此类体在单个占位符下嵌入多行重载集。 \texttt{gather\_docstrs} 为 \texttt{tensorplay.\_C.\_add\_docstr} 安装 \texttt{Mock} 并单独执行 \texttt{\_docs.py} ，然后 \texttt{\_add\_docstr\_to\_hint} 将 \texttt{...} 之后的文档字符串与 \texttt{textwrap.indent} 和 \texttt{fmt: skip} 防护进行拼接。

\begin{lstlisting}[caption={template substitution.}, label=lst:pyisubst, style=pystyle]
def _substitute(template: str, env: dict[str, str]) -> str:  # gen_pyi.py:500
    def replace(m):
        indent, key = m.group(1), m.group(2)
        if indent is None:  # mid-line hole: verbatim
            return env[key]
        content = textwrap.indent(env[key], indent)  # block hole
        return "\n".join(line.rstrip() for line in content.splitlines())
    return _SUBST_RE.sub(replace, template)
\end{lstlisting}

\subsection{CPython 快速路径：（944 行）}

\texttt{gen\_python\_c.py} 拥有绕过 \texttt{pybind11} 过载调度的 \texttt{METH\_FASTCALL | METH\_KEYWORDS} 层。 \texttt{\_BRIDGE} 将每个 C++ 参数类型映射到 \texttt{CPythonBridge.h} 解包器（\texttt{tpx\_py\_tensor\_cref}、\texttt{tpx\_py\_int64}、\texttt{tpx\_py\_intlist} 等）； \texttt{\_KIND\_CONST} 将其映射到 \texttt{TPK\_INT} / \texttt{TPK\_TENSOR} 类型字节以进行急切类型检查。 \texttt{\_emit\_op} 每次重载发出一个本机条目：它构建 \texttt{kwlist}，通过 \texttt{tpx\_py\_parse\_into} 解析位置和关键字参数，验证所需参数，对静态 \texttt{tpx\_kinds} 表执行类型检查，通过 \texttt{\_BRIDGE} 解包，释放 \texttt{tpx::ops::} 调用周围的 GIL，并通过 \texttt{\_pack\_expr} 打包结果。多重重载名称获得一个组调度程序，该组调度程序按声明顺序尝试候选者，仅在 \texttt{std::invalid\_argument} （参数形状不匹配）上失败，而内核故障会立即转换。该模块通过 \texttt{register\_generated\_cpython\_functions} 和 \texttt{register\_generated\_cpython\_methods} 安装表，跳过任何已手动绑定的名称，以便工厂数据类型/设备解析语义保持权威。

\begin{lstlisting}[caption={per-overload skeleton (abridged, one \texttt{add} overload).}, label=lst:cpythonop, style=cppstyle]
static PyObject* pyop_add_function_ov0(PyObject*, PyObject* const* args,
                                       Py_ssize_t nargs, PyObject* kwnames) {
    static const char* kwlist[] = {"self", "other", "alpha", nullptr};
    PyObject* slots[3];
    tpx_py_parse_into(args, nargs, kwnames, kwlist, 3, "add.Tensor", slots);
    static const unsigned char tpx_kinds[] = {TPK_TENSOR, TPK_TENSOR, TPK_NUMBER};
    tpx_py_check_types(slots, 3, "add.Tensor", kwlist, tpx_kinds, 2);
    auto&& s_self  = tpx_py_tensor_cref(slots[0]);
    auto&& s_other = tpx_py_tensor_cref(slots[1]);
    auto&& s_alpha = tpx_py_scalar(slots[2]);
    tensorplay::python::tpx_prof_capture_site();
    auto r = [&](){ tpx_py_GilRelease _gil;
                    return static_cast<Tensor(*)(const Tensor&,const Tensor&,Scalar)>
                           (tensorplay::tpx::ops::add)(s_self,s_other,s_alpha); }();
    return tpx_py_wrap(r);
}
\end{lstlisting}

\subsection{端到端流程}

图~\ref{fig:codegen_contract} 显示了从 YAML 到调度的管道。

\begin{figure}[H]
\centering
\begin{tikzpicture}[node distance=1.1cm, >=Stealth, font=\small]
\tikzset{box/.style={draw, rounded corners=3pt, thick, align=center, minimum width=2.6cm, minimum height=1.15cm, fill=codebg},
         arr/.style={->, thick, color=tpblue}}
\node[box, fill=tporange!12] (yaml) {\texttt{native\_functions.yaml}\\479 names / 2\,809 records\\+ \texttt{derivatives.yaml}\\187 formulas};
\node[box, right=1.6cm of yaml, fill=tpblue!10] (model) {\\ \texttt{parse\_schema} via local\\schema model + \texttt{NativeFunction}};
\node[box, right=1.6cm of model, fill=tpgreen!10] (ctx) {\\ \texttt{CodegenContext}\\ \texttt{DEFAULT\_TARGETS} (11)};
\node[box, below=1.25cm of ctx, fill=tppurple!10] (gen) {\\\\\\ 12 generators};
\node[box, left=1.6cm of gen, fill=tpblue!10] (stage) {\\ \texttt{generated}\\10 headers + \texttt{functional.py}};
\node[box, left=1.6cm of stage] (compile) {\\ \texttt{\#include TensorGenerated.h}\\inside \texttt{class Tensor}};
\node[box, below=1.25cm of compile, fill=tporange!12] (dispatch) {Link \& dispatch\\ \texttt{detail::redispatch\_*}\\ \texttt{DispatchStub::call}\\ \texttt{ GpuTimerPair}};
\draw[arr] (yaml) -- (model);
\draw[arr] (model) -- (ctx);
\draw[arr] (ctx) -- (gen);
\draw[arr] (gen) -- (stage);
\draw[arr] (stage) -- (compile);
\draw[arr] (compile) -- (dispatch);
\draw[arr, dashed, color=tpgreen] (yaml) -- ++(0,-3.05) -| (dispatch) node[pos=0.22, below, font=\scriptsize, text=gray] {re-parse on every build};
\end{tikzpicture}
\caption{合同$\rightarrow$生成$\rightarrow$调度。 YAML 的一次解析可提供 11 个生成器；保证在任何编译之前重新生成；包含 at 使生成的成员成为普通 \texttt{Tensor} 成员。}
\label{fig:codegen_contract}
\end{figure}

\subsection{工件库存和尺寸}

干净生成后，十个阶段输出占用磁盘上的以下预算：

\begin{table}[H]
\centering
\caption{生成的工件大小（CPU 版本，2026-09-02）。 Python 包输出暂存于源代码树中。}
\label{tab:codegen_sizes}
\small
\begin{tabular}{lrrl}
\toprule
人工制品 & 线路 & 字节 & Note \\
\midrule
\texttt{TensorGenerated.h} & 6\,900 & 355\,k & 479 个不同的名称，在重载键上进行重复数据删除 \\
\texttt{TensorGenerated.cpp} & 12\,200 & 9\,799\,k & 两遍：重新调度+成员 \\
\texttt{TensorRedispatchGenerated.h} & 4\,200 & 465\,k & 一份声明每（操作，变体） \\
\texttt{TPXOpsGenerated.h} & 5\,100 & 592\,k & 标题中的内联简单转发器 \\
\texttt{TPXOpsGenerated.cpp} & 4\,800 & 453\,k & 具有图形逻辑的外线包装器 \\
\texttt{TPXAutogradRegistration.cpp} & 1\,900 & 200\,k & 静态构造函数，187 个内核 $\times$~3 个键 \\
\texttt{AutogradNodesGenerated.h} & 3\,200 & 286\,k & 148 个节点结构 \\
\texttt{AutocastGenerated.cpp} & 1\,100 & 111\,k & 自动铸造内核 \\
\texttt{TensorCPythonGenerated.h} & 98\,000 & 9\,304\,k & 每次过载 \texttt{METH\_FASTCALL} \\
\texttt{TensorBindingsGenerated.h} & 12 & $<1$\,k & 后备 pybind11 表面 \\
\texttt{functional.py} & 950 & 48\,k & 上演至 \texttt{pkg\_out} \\
\texttt{*.pyi} & 2\,400 & 96\,k & 通过 \texttt{gen\_pyi.py} 的三个存根 \\
\bottomrule
\end{tabular}
\end{table}

大小排序是故意的：重新调度层和 CPython FASTCALL 层占主导地位，因为它们显式枚举每个重载；手写的调度程序表（超过 13 个键的 \texttt{DispatchTable}）保持较小且不变。 \texttt{TensorGenerated.cpp} 和 \texttt{TensorCPythonGenerated.h} 中操作员数量的增长呈线性；添加一个新模式正好添加一个重新调度渠道、一个成员和一个 CPython 条目，通过比较 YAML 编辑前后的暂存树进行验证。

\subsection{新鲜度和失效模式}

过时的工件并不微妙：取决于每个 \texttt{*.py}，因此触摸任何生成器或模式解析器都会强制在下一个 \texttt{ninja -C build} 链接之前重新生成。 \texttt{GENERATED TRUE} 属性告诉 CMake 删除 \texttt{cmake --build --target clean} 上的输出，并且是 \texttt{.gitignore} 的，因此不会意外提交过时的标头。失败是快速失败的：\texttt{\_BRIDGE} (\texttt{\_require\_bridge}) 中的未知 C++ 类型会引发 \texttt{SystemExit} 和 \texttt{native CPython bridge has no converter for C++ type}，\texttt{\_default\_pyobject} 中未处理的默认值会引发相同的情况，调度投影不匹配会打印预期的和实际的内核映射。最常见的用户错误是手动编辑 \texttt{TensorGenerated.h}；文件头 \texttt{main.py -- DO NOT EDIT} 是一个契约：下一个版本会覆盖它，并且 \texttt{git diff} 显示漂移。


